\textbf{Job-Related Health and Safety}
\par
As a software engineer, it is crucial to
prioritize your health and safety while working.
Ensure your workstation is ergonomically set up
to prevent strain and injury, including an adjustable
chair, properly positioned monitor, and correctly
placed keyboard and mouse.
\par
\bigskip
Regular breaks to stand, stretch, and move around
are essential to avoid prolonged sitting. Protect
your eye health by using the 20-20-20 rule, adjusting
screen brightness, and employing anti-glare filters.
To prevent repetitive strain injuries, practice good
typing habits and consider using ergonomic keyboard
and mouse devices.
\par
\bigskip
Maintaining a healthy work-life balance, setting
boundaries, and taking regular breaks can help
manage stress and prevent burnout. Incorporating
physical activity, such as short walks or stretching
exercises, counteracts the sedentary nature of desk work.
\par
\bigskip
Keep your workspace clean and organized to prevent
accidents, ensure electrical equipment is well-maintained,
and manage cables to avoid tripping hazards.
Adequate lighting, preferably natural or full-spectrum
lighting, reduces eye strain and improves focus.
\par
\bigskip
By following these guidelines, you can create a more
comfortable and productive work environment, reduce
the risk of injury, and maintain overall well-being
as a software engineer.
\par
\bigskip
\textbf{Food and Drink Guidelines}
\par
Students will not be allowed to have open drink
containers in the classroom. Students are encouraged
to eat in the break area or outside areas.
\textbf{Drug, Alcohol, and Tobacco Free Campus}
\par
Dixie Technical College is a drug, alcohol,
tobacco, and vape free campus. See College
policies for more information.
\textbf{Cell Phone Guidelines}
\par
Students may use cell phones while on campus.
We ask that students put cell phones on silent
during class times. If the use of cell phones
or headphones causes disruptions or loss in
productivity, the student may become subject
to disciplinary action. We expect students to
use their cell phones responsibly for learning only.
